% Cryptographic Settlement Binding for Confidential On-Chain State — Harpo (v3)
% Compila com: pdflatex instant-settlement-v3.tex (2x para refs/citacoes)
% Dependencias: pacotes padrao (article, amsmath, amssymb, amsthm, hyperref, booktabs, xcolor, geometry).
% Bibliografia embutida via thebibliography (sem BibTeX), pela mesma razao de
% reprodutibilidade adotada na v2: nenhuma dependencia exotica.
\ifdefined\XeTeXversion\else\pdfoutput=1\fi
\documentclass[10pt,twocolumn]{article}

\usepackage[utf8]{inputenc}
\usepackage[T1]{fontenc}
\usepackage{lmodern}
\usepackage{amsmath,amssymb,amsthm}
\usepackage{booktabs}
\usepackage{xcolor}
\usepackage[margin=0.75in]{geometry}
\setlength{\columnsep}{0.25in}
\usepackage{microtype}
\usepackage[colorlinks=true,linkcolor=blue,citecolor=blue,urlcolor=blue,breaklinks=true,
  pdftitle={Cryptographic Settlement Binding for Confidential On-Chain State},
  pdfauthor={Juliano Pereira Sales, Marco Tulio Rocha Nascimento},
  pdfsubject={Zero-knowledge proofs, confidential settlement, instant payments},
  pdfkeywords={zero-knowledge proofs, confidential transactions, settlement binding, payment attestation, external-state binding, EVM, Groth16, Poseidon, instant payments}
]{hyperref}
\usepackage{xurl}
\usepackage{longtable}
\setlength{\emergencystretch}{1.5em}

\newtheorem{property}{Property}
\newtheorem{theorem}{Theorem}

\newcommand{\Poseidon}{\ensuremath{\mathsf{Poseidon}}}
\newcommand{\eeid}{\ensuremath{\mathsf{e2eId}}}
\newcommand{\code}[1]{\texttt{#1}}
\newcommand{\amount}{\ensuremath{\text{amount}}}
\newcommand{\txid}{\ensuremath{\text{txid}}}
\newcommand{\party}{\ensuremath{\text{party}}}
\newcommand{\salt}{\ensuremath{\text{salt}}}
\newcommand{\pc}{\ensuremath{\mathit{pc}}}
\newcommand{\Setup}{\ensuremath{\mathsf{Setup}}}
\newcommand{\SetupZK}{\ensuremath{\mathsf{Setup\_ZK}}}
\newcommand{\SetupDomain}{\ensuremath{\mathsf{Setup\_Domain}}}
\newcommand{\Commit}{\ensuremath{\mathsf{Commit}}}
\newcommand{\Prove}{\ensuremath{\mathsf{Prove}}}
\newcommand{\Verify}{\ensuremath{\mathsf{Verify}}}
\newcommand{\Attest}{\ensuremath{\mathsf{Attest}}}
\newcommand{\Settle}{\ensuremath{\mathsf{Settle}}}
\newcommand{\Locked}{\ensuremath{\mathsf{Locked}}}
\newcommand{\Finalized}{\ensuremath{\mathsf{Finalized}}}
\newcommand{\Irrevocable}{\ensuremath{\mathsf{Irrevocable}}}

\title{\textbf{Cryptographic Settlement Binding}\\
for Confidential On-Chain State}

\author{%
  \begin{tabular}[t]{c}
  Juliano Pereira Sales \\
  \small Harpo-ZK, Brazil \\
  \small \href{mailto:tzdesing@gmail.com}{tzdesing@gmail.com}
  \end{tabular}
  \and
  \begin{tabular}[t]{c}
  Marco T\'ulio Rocha Nascimento \\
  \small Harpo-ZK, Brazil \\
  \small \href{mailto:marcotulio.nascimento@gmail.com}{marcotulio.nascimento@gmail.com}
  \end{tabular}
}

\date{\today}

\makeatletter
\newenvironment{widetitleabstract}{\begin{@twocolumnfalse}}{\end{@twocolumnfalse}}
\makeatother

\begin{document}

\twocolumn[
\begin{widetitleabstract}
\maketitle
\begin{abstract}
A recurring problem in confidential blockchain applications is the need to
establish a cryptographic correspondence between a private on-chain
obligation created before an external settlement event and an externally
observed settlement event that becomes available only later. A
zero-knowledge proof can establish the internal consistency of the
on-chain state, but it does not by itself establish that an external
payment actually occurred. Conversely, an institution may authenticate a
statement about an external payment without revealing the payment details
on-chain.

We present \textbf{cryptographic settlement binding}, a construction that
connects these two states while keeping the settlement amount, payment
identifier, and counterparty identifier confidential. The construction is
two-phase. Before settlement, a commitment $c_0$ binds the agreed amount
and operation identifier to an on-chain obligation. After settlement, a
second commitment $c$ binds the same amount and operation identifier to
the settlement-side payment identifier. A domain-keyed tag $\nu$ provides
a deterministic, single-use reference to the settlement identifier
without revealing it, while an on-chain domain anchor $K_D$ binds the tag
to the registered domain key. A registered payment service provider (PSP)
signs an EIP-712 receipt over the commitments rather than over plaintext
payment data.

We formalize the resulting properties under Groth16 knowledge soundness,
Poseidon and Keccak collision resistance, and ECDSA unforgeability. The
construction proves cryptographic consistency and authenticates that a
registered PSP signed the accepted statement; it does \textbf{not}
cryptographically establish that the underlying payment rail was
truthful. That remaining correspondence is stated explicitly as an
institutional trust assumption. This separation prevents authentication
of an institutional statement from being confused with independent
evidence of external-world truth.

A reference implementation uses Circom, Groth16 and Solidity and is
exercised on a public EVM testnet. The settlement circuit contains 1,251
R1CS constraints at the reported optimization level, and the reference
\code{finalize} transaction costs 353,062 gas. A concrete Pix
instantiation is provided as an application profile; the construction
itself is independent of Pix.

\medskip
\noindent\textbf{Keywords:} zero-knowledge proofs, confidential
transactions, settlement binding, payment attestation, external-state
binding, EVM, Groth16, Poseidon, instant payments
\end{abstract}
\end{widetitleabstract}
\vspace{1em}
]

\section{Introduction}

A confidential state transition on a public blockchain can establish that
an on-chain operation is internally consistent without revealing the
private data on which that consistency depends. It can, for example,
prove knowledge of an amount, ownership relation, or transaction witness
while exposing only commitments and zero-knowledge proofs. This
establishes a property of the blockchain state itself.

A different problem arises when the state transition depends on an event
that occurs outside the blockchain. Consider an operation in which an
asset or obligation is committed on-chain before an external payment is
made. Afterward, an authorized institution observes a settlement event
and obtains information identifying that event. The system must then
establish that the observed settlement corresponds to the particular
confidential obligation that was committed earlier. The difficulty is not
merely to prove that the external event exists, nor merely to prove that
a private on-chain transition is well formed. It is to establish a
cryptographic correspondence between the two events while keeping the
sensitive attributes of both confidential.

We call this problem \textbf{cryptographic settlement binding}.

Informally, a settlement-binding protocol allows a verifier to establish
that an externally observed settlement event is consistent with a
confidential state transition previously committed on-chain, while
revealing only the information required by the application. The temporal
order is essential. Before settlement, the system can commit to the terms
of an obligation, but the identifier of the eventual settlement event is
not yet available as evidence of that event. After settlement, an
authorized observer can identify and authenticate the external event, but
the system must still demonstrate that it corresponds to the state
committed earlier.

A straightforward design could reveal the amount and external payment
identifier to the blockchain and have a contract compare them with the
original obligation. Such a design would establish the desired
correspondence, but would defeat the confidentiality requirement.
Alternatively, a zero-knowledge proof could hide the values while proving
that a private payment identifier is related to an on-chain commitment.
This alone, however, does not identify who is authorized to make the
external settlement statement, nor does it distinguish cryptographic
consistency from the truthfulness of that statement.

Our construction separates these concerns.

The first component is a \textbf{two-phase commitment}. Before
settlement, an operation commits to its agreed terms as
\[
c_0 = \Poseidon(\amount, \salt_0, \txid).
\]
The commitment is registered on-chain before the external settlement is
finalized and can therefore serve as an anchor for an asset escrow or
another dependent state transition. After settlement, a second commitment
is formed as
\[
c = \Poseidon(\amount, \salt, \eeid, \txid),
\]
where \eeid\ denotes the settlement-side identifier supplied by the
concrete payment rail. The settlement proof demonstrates knowledge of a
common \amount\ and \txid\ underlying both commitments. Thus the external
settlement identifier is introduced only in the second phase, while the
relationship to the previously committed obligation is established in
zero knowledge.

The second component is a \textbf{domain-keyed settlement tag},
\[
\nu = \Poseidon(\eeid, k_D), \qquad K_D = \Poseidon(k_D).
\]
The tag gives the on-chain protocol a deterministic identifier for the
settlement under a particular domain without exposing the underlying
settlement identifier. In the reference construction, the tag is consumed
at most once by a deployed settlement instance. This provides a
cryptographically checkable single-use property while preserving
confidentiality of the external identifier. The domain key also makes the
linkage selective: an observer that does not possess $k_D$ cannot simply
recompute tags for arbitrary settlement identifiers. The resulting
property is not global anonymity or unlinkability; those properties are
outside the scope of the construction.

The third component is an explicit \textbf{institutional trust boundary}.
A registered settlement observer signs an authenticated receipt
containing the pre-settlement commitment and the public settlement-binding
values. In the reference construction, the receipt is represented using
EIP-712 and includes an expiration condition. The zero-knowledge proof
establishes consistency between the commitments and the settlement
identifier; the signature establishes that the statement was issued by a
registered attestor. These are cryptographic statements about consistency
and authentication.

They are not, by themselves, statements that the external payment rail
was truthful.

More precisely, acceptance establishes that the accepted proof is
consistent with the pre-settlement commitment and that the corresponding
receipt was authenticated by a registered attestor. The construction
therefore makes the remaining institutional assumption explicit: the
authorized attestor is assumed to report settlement correctly. This
assumption is not hidden inside the zero-knowledge relation and is not
presented as a consequence of cryptography. Removing or weakening it is
treated separately as an open problem.

This separation is important because an authenticated statement and a
truthful statement are different security objects. A cryptographic
protocol can prevent an adversary from forging a statement under an
honest attestor's key; it cannot, without an additional source of trusted
evidence, determine whether the institution itself reported a false
settlement. The construction consequently reduces the cryptographic
problem to authentication and consistency while exposing the
institutional assumption as an explicit boundary of the security model.

The paper makes three main contributions.

\textbf{First, we define and construct cryptographic settlement binding.}
The construction introduces a two-phase commitment relation in which a
confidential obligation is committed before settlement and a subsequent
zero-knowledge proof binds the resulting external settlement identifier
to the same confidential terms. The construction is expressed
independently of a particular payment rail.

\textbf{Second, we introduce a domain-keyed settlement tag for
confidential single-use identification.} The tag deterministically
represents the external settlement identifier under a domain secret,
while its public representation does not disclose that identifier. The
security analysis characterizes its binding and single-use properties
under the stated cryptographic assumptions and deployment scope.

\textbf{Third, we make the institutional trust boundary explicit and
characterize what the cryptographic protocol does and does not
establish.} In particular, the security analysis distinguishes
authentication of a settlement statement from the truth of the underlying
external event. This separation prevents the protocol's cryptographic
guarantees from being conflated with assumptions about the behavior of
the institution providing the attestation.

The remainder of the paper formalizes these claims. Section~\ref{sec:model}
defines the settlement-binding model, its participants, state, and
security goal. Section~\ref{sec:construction} gives the construction and
its protocol algorithms. Section~\ref{sec:security} analyzes commitment
binding, two-phase consistency, attestation authenticity, and single-use
settlement. Section~\ref{sec:privacy} characterizes the resulting
confidentiality properties and their limits. Section~\ref{sec:pix}
instantiates the construction for Pix, while Section~\ref{sec:predicates}
describes policy predicates that can consume the resulting confidential
commitments. Section~\ref{sec:impl} evaluates the reference
implementation, including circuit size, gas consumption, and
well-determinedness analysis. Section~\ref{sec:related} positions the
construction relative to external-payment proofs and confidential EVM
systems, and Section~\ref{sec:limitations} states the remaining
limitations and open problems.

The construction is intentionally narrower than a complete payment
system. It does not establish central-bank or payment-rail truth without
an additional trusted source, does not provide transaction unlinkability,
and does not claim economic finality beyond the state-transition policy
of the deployment. Pix is used only as a concrete instantiation of the
generic construction in Section~\ref{sec:pix}; the cryptographic
construction and its security properties do not depend on Pix-specific
semantics.

\section{Settlement-Binding Model}
\label{sec:model}

\subsection{Problem Definition}

We consider an application in which a confidential on-chain operation
depends on an external settlement event.

Let an operation $O$ have agreed private terms $T=(\amount,\txid)$, where
\amount\ is the agreed value and \txid\ is an application-level operation
identifier. Before the external settlement occurs, the application must
be able to create an on-chain state that commits to these terms.

Later, an external payment or settlement rail produces an identifier
\eeid, observable by an authorized institution. The identifier \eeid\ is
not itself assumed to prove that settlement occurred: an identifier may
exist for an operation that is subsequently rejected or otherwise fails
to settle. The external event therefore requires an authenticated
statement from an institution that is authorized to observe the rail.

The problem addressed by this paper is to establish a binding between the
two states: (1) a confidential commitment to the agreed terms created
before settlement; and (2) a confidential commitment incorporating the
settlement-side identifier created after settlement.

We call this correspondence \textbf{cryptographic settlement binding}.

The construction is intentionally independent of a particular payment
rail. A rail-specific deployment supplies an identifier \eeid, a
canonical representation of the relevant external fields, and an
authorized attestor capable of evaluating the required commitment. The
core construction does not require the settlement identifier itself to
become public.

\subsection{Participants and Trust Boundary}

The model contains four logical roles.

\textbf{Obligation issuer.} Creates the pre-settlement operation and its
commitment. The issuer may be a smart contract, an application, or an
escrow component.

\textbf{Prover.} Holds the private witness required to demonstrate that
the pre- and post-settlement commitments refer to the same operation.

\textbf{Attestor.} An authorized institution that observes the external
settlement system and signs a statement about the observed event.

\textbf{Verifier.} The on-chain contract that accepts or rejects the
settlement transition.

The construction separates two claims that are often conflated.

First, the verifier can establish a cryptographic statement:
\begin{quote}
\small
\emph{same confidential witness} $\Rightarrow$ \emph{consistent pre- and
post-settlement commitments.}
\end{quote}
Second, the verifier can establish an authentication statement:
\begin{quote}
\small
\emph{registered attestor} $\Rightarrow$ \emph{the accepted receipt was
signed by an authorized key.}
\end{quote}
Neither statement alone establishes that the external settlement system
actually behaved as claimed. That final correspondence is an
institutional assumption about the attestor.

Accordingly, the principal trust boundary is:
\begin{quote}
\small
\fbox{\begin{minipage}{0.9\linewidth}
\centering
cryptographic consistency $+$ attestor authentication $\neq$
independent rail truth
\end{minipage}}
\end{quote}
The distinction is fundamental to the security claims of the
construction.

\subsection{Commitments and Domain Separation}
\label{sec:model-commitments}

Let $H$ be a collision-resistant hash function suitable for use inside
the selected zero-knowledge proof system.

For an operation with \amount, operation identifier \txid, settlement
identifier \eeid, and fresh randomness $\salt_0,\salt$, define
conceptually:
\[
\begin{aligned}
c_0 &= H(\amount,\salt_0,\txid) \\
c &= H(\amount,\salt,\eeid,\txid).
\end{aligned}
\]
The first commitment represents the obligation before settlement. The
second represents the corresponding state after settlement.

A domain $D$ is associated with a secret key $k_D$ and a public anchor
\[
K_D = H(k_D).
\]
The settlement-side identifier is represented by a domain-keyed tag
\[
\nu = H(\eeid,k_D).
\]
The tag provides a public handle that does not reveal \eeid\ to an
observer that does not possess the domain key.

A counterparty commitment may similarly be represented as
\[
\pc = H(\party,\salt_p),
\]
where \party\ is the relevant counterparty identifier and $\salt_p$ is
fresh randomness.

These equations define the abstract relation that the construction must
prove. Their concrete circuit realization is given in
Section~\ref{sec:construction}.

\subsection{Settlement-Bound Acceptance}

An operation is \textbf{cryptographically settlement-bound} when an
accepting execution establishes all of the following:
\begin{enumerate}
\item a valid pre-settlement commitment $c_0$ was registered for the
  operation;
\item a valid zero-knowledge witness connects $c_0$ and $c$ through the
  same \amount\ and operation identifier;
\item the settlement tag $\nu$ is derived from the settlement-side
  identifier \eeid\ under the registered domain key;
\item the accepted counterparty commitment is part of the same witness;
\item a registered attestor authenticated the same $(c_0,\nu,\pc)$ tuple
  accepted by the verifier; and
\item the tag cannot be consumed twice within the deployed settlement
  instance.
\end{enumerate}
The resulting statement is therefore stronger than either a standalone
zero-knowledge transfer proof or a standalone payment signature: it binds
an external settlement observation to a previously registered
confidential on-chain obligation.

\subsection{Security Properties}

The formal analysis in Section~\ref{sec:security} considers four
principal properties.

\textbf{Commitment consistency.} An accepting proof cannot substitute
different values or operation identifiers between the pre- and
post-settlement commitments, except with negligible probability under the
stated cryptographic assumptions.

\textbf{Settlement-identifier binding.} An accepting proof binds the
public settlement tag to the settlement identifier contained in its
witness under the registered domain key.

\textbf{Attestation binding.} An accepting settlement is associated with
a receipt signed by a registered attestor over the same public
commitments accepted by the verifier.

\textbf{Instance-scoped single use.} A settlement tag accepted by one
deployed settlement instance cannot be consumed again by that instance.

The properties deliberately do not include a claim that the external
payment rail itself is truthful. That statement belongs to the
institutional assumption described next.

\subsection{Institutional Assumption}

Let $I_1$ denote the trusted-attestor assumption.

A registered attestor signs $\mathit{settled}=\mathit{true}$ only for an
external payment that it observed as settled and credited, and maintains
an authoritative internal mapping between the operation identifier,
settlement identifier, amount, and relevant counterparty information.

Under $I_1$, the attestation provides an authenticated bridge from the
external event to the cryptographic construction.

Without $I_1$, the protocol still authenticates that a registered key
signed the statement, but it cannot distinguish a truthful statement from
a deliberately false one produced by a malicious or compromised attestor.

This distinction is preserved throughout the remainder of the paper.

\subsection{What the Model Does Not Claim}

The model does not assume that the external rail exposes a
cryptographically verifiable proof to the public EVM. It also does not
provide economic finality, anonymity, or unlinkability merely from the
use of zero-knowledge proofs.

In particular, confidentiality of \amount, \eeid, and \party\ from
ordinary chain observers is distinct from unlinkability between
operations. Holders of the domain key can intentionally use $\nu$ as a
selective-linkability mechanism.

The concrete construction implementing this model is specified in
Section~\ref{sec:construction}.

\section{Settlement-Binding Construction}
\label{sec:construction}

This section instantiates the abstract model of Section~\ref{sec:model}
as a concrete EVM-compatible construction using Poseidon commitments
\cite{poseidon}, a Groth16 proof \cite{groth16}, a domain-keyed settlement
tag, and an EIP-712 attestation receipt. The construction is represented
by the tuple
\[
\Pi = (\Setup, \Commit, \Prove, \Verify, \Attest, \Settle).
\]
The algorithms are not intended to introduce a new zero-knowledge proof
system. They specify how standard commitments, a Groth16 proof, an
institutional attestation, and an on-chain state transition are composed
to realize the settlement-binding relation defined in
Section~\ref{sec:model}.

\subsection{Setup}

\Setup\ is the composition of two independent sub-algorithms,
\[
\Setup \;\rightarrow\; (\SetupZK, \SetupDomain),
\]
which are not run together and do not share randomness or inputs.

\paragraph{\SetupZK.}
\[
\SetupZK(1^\lambda) \;\rightarrow\; (pp,vk)
\]
generates the Groth16 proving and verification parameters for the
settlement circuit. This step requires a trusted-setup ceremony and is
run once per circuit, independently of any settlement domain. The same
$(pp,vk)$ pair is reused by every domain that deploys the same circuit.

\paragraph{\SetupDomain.}
\[
\SetupDomain(1^\lambda, D) \;\rightarrow\; (k_D, K_D)
\]
samples a fresh domain secret $k_D$ for a settlement domain $D$ and
computes the public domain anchor
\[
K_D = \Poseidon(k_D).
\]
Creating a domain never requires a new Groth16 ceremony. The deployed
verifier stores $K_D$ as the domain anchor; rotating $k_D$ creates a new
domain rather than mutating the existing anchor. A deployment may create
many domains against a single $(pp,vk)$ produced once by \SetupZK.

\subsection{Commit}

For agreed private terms $T=(\amount,\txid)$, the commitment algorithm
samples fresh $\salt_0$ and computes
\[
\begin{aligned}
c_0 &= \Commit_0(\amount,\txid;\salt_0) \\
    &= \Poseidon(\amount,\salt_0,\txid).
\end{aligned}
\]
The application registers $c_0$ with the on-chain settlement instance and
locks the corresponding operation. At this point the settlement-side
identifier \eeid\ is not required by the on-chain state: the operation
establishes its cryptographic identity before the external event occurs.

After settlement, the prover samples fresh \salt\ and computes
\[
\begin{aligned}
c &= \Commit_1(\amount,\eeid,\txid;\salt) \\
  &= \Poseidon(\amount,\salt,\eeid,\txid),
\end{aligned}
\]
where \eeid\ is the settlement-side identifier obtained from the external
payment rail.

The two commitments are intentionally distinct: $c_0$ identifies the
obligation before settlement, while $c$ incorporates the external
settlement identifier. The same private \amount\ and \txid\ are reused
across both, rather than independently supplied, which is what gives the
construction its binding property (made explicit in
Section~\ref{sec:invariant}).

\subsection{Prove}
\label{sec:prove}

The prover constructs the witness
\[
w = (\amount,\salt_0,\salt,\eeid,\txid,k_D,\party,\salt_p)
\]
and computes
\[
\begin{aligned}
\nu &= \Poseidon(\eeid,k_D), \\
\pc &= \Poseidon(\party,\salt_p).
\end{aligned}
\]
The public statement is
\[
y = (c_0,c,\nu,K_D,\pc).
\]
The prover then computes
\[
\pi \leftarrow \Prove(pp,w;y).
\]
The circuit accepts the witness only if
\begin{align}
c_0 &= \Poseidon(\amount,\salt_0,\txid), \label{eq:c0}\\
c &= \Poseidon(\amount,\salt,\eeid,\txid), \label{eq:c}\\
\nu &= \Poseidon(\eeid,k_D), \label{eq:nu}\\
K_D &= \Poseidon(k_D), \label{eq:KD}\\
\pc &= \Poseidon(\party,\salt_p). \label{eq:pc}
\end{align}
The same private \amount\ and \txid\ are structurally constrained in both
phases: a valid proof does not merely show that two independently
supplied hashes are valid, but that both commitments were derived from
the same private amount and operation identifier.

The settlement tag $\nu$ is intentionally scoped to the domain key $k_D$
rather than computed directly from \eeid. This prevents an observer who
already knows a candidate \eeid\ from testing it directly against
on-chain tags without knowledge of $k_D$, and it makes the single-use
property enforced on $\nu$ (Section~\ref{sec:settle}) instance-scoped
rather than globally unique across independently deployed domains.

The counterparty commitment \pc\ is likewise not published in plaintext;
it is cross-checked against the PSP's own commitment at settlement time
(Section~\ref{sec:attest}), which is what makes downstream predicates
over \pc\ meaningful rather than decorative --- without that cross-check,
a prover could generate a valid settlement proof while associating it
with an unrelated counterparty commitment.

\subsection{Verify}

The proof verifier evaluates
\[
\Verify(vk,y,\pi) \in \{0,1\}.
\]
Acceptance establishes only the cryptographic relation represented by the
circuit. In particular, successful verification does not independently
establish that the external settlement identified by \eeid\ actually
occurred. The latter statement enters through \Attest.

\subsection{Attest}
\label{sec:attest}

A registered payment institution evaluates the external event and
constructs the typed receipt
\[
R = (c_0,\nu,\pc,\mathit{settled},\mathit{deadline}).
\]
The attestor signs the receipt using EIP-712:
\[
\sigma \leftarrow \mathsf{Sign}_{sk_P}(R).
\]
The attestation therefore authenticates the public commitments rather
than exposing the amount, settlement identifier, or counterparty
identifier.

The attestor is assumed to derive $c_0$ from the amount it actually
observed for the external settlement --- independently recomputing
\[
c_0' = \Poseidon(\mathit{amount\_settled},\salt_0,\txid)
\]
and signing only when $c_0'$ matches the operation's registered $c_0$.
This gives the construction its \textbf{commitment-carried
value-consistency} property: the attestor authenticates a commitment
derived from the observed amount without the amount itself ever
appearing in the receipt.

The receipt also carries an expiration deadline; a stale receipt cannot
be used to finalize an operation after its authorization period.

\subsection{Settle}
\label{sec:settle}

The on-chain settlement algorithm evaluates
\[
\Settle(c_0,\nu,\pc,\pi,R,\sigma).
\]
It accepts only if:
\begin{enumerate}
\item the operation is currently \Locked;
\item $\sigma$ authenticates a currently registered attestor;
\item $\mathit{settled}=\mathit{true}$;
\item the receipt has not expired;
\item the receipt's $c_0$ equals the locked commitment;
\item the receipt's \pc\ equals the proof's public \pc;
\item the proof's $K_D$ equals the registered domain anchor;
\item $\Verify(vk,y,\pi)=1$; and
\item $\nu$ has not already been consumed.
\end{enumerate}
Upon acceptance, the contract records the settlement result and consumes
$\nu$, transitioning the operation to \Finalized.

The ordering of the checks is also relevant operationally: inexpensive
receipt and state checks (1--7) can reject malformed submissions before
the pairing-based proof verification (8) is executed.

\subsection{Construction Invariant}
\label{sec:invariant}

The central invariant of the construction is that
\[
\begin{aligned}
c_0 &= \Poseidon(\amount,\salt_0,\txid) \\
c &= \Poseidon(\amount,\salt,\eeid,\txid)
\end{aligned}
\]
share the private values \amount\ and \txid. Because the same private
amount and txid occur in both relations, an accepting proof establishes
that the pre-settlement and post-settlement commitments refer to the
same agreed value and operation identifier, subject to the cryptographic
assumptions stated in Section~\ref{sec:security}.

The additional relations
\[
\begin{aligned}
\nu &= \Poseidon(\eeid,k_D) \\
K_D &= \Poseidon(k_D)
\end{aligned}
\]
bind the public settlement tag to the registered domain.

Finally, the attestation relation binds the accepted public tuple
$(c_0,\nu,\pc)$ to a registered institutional key.

The security consequences of these relations are proved in
Section~\ref{sec:security}.

\subsection{Reversal and Liveness}
\label{sec:reversal}

A valid attestation does not necessarily imply economic irreversibility
of the external settlement.

The reference deployment therefore distinguishes \Finalized\ from
\Irrevocable.

Within the configured dispute window, the authorized attestor may
register a reversal. A reversed operation cannot become releasable merely
because the nominal window subsequently expires.

Separately, an operation that remains locked without successful
settlement can expire through the protocol timeout. This prevents an
unavailable or non-cooperating attestor from permanently locking the
pre-settlement obligation.

The timeout is a liveness mechanism; it is not evidence about the
external payment rail.

\subsection{Protocol Summary}

The resulting sequence realizes $\Pi$ end to end:
\begin{enumerate}
\item run \SetupDomain\ to register the domain anchor $K_D$ (\SetupZK\
  having already produced $(pp,vk)$ once for the circuit);
\item run \Commit\ to compute $c_0$ from the agreed amount and operation
  identifier, and lock the operation against $c_0$;
\item observe the external settlement and obtain the settlement-side
  identifier \eeid;
\item run \Commit\ to construct $c$, and run \Prove\ to construct $\nu$,
  \pc, and the proof $\pi$;
\item obtain the PSP receipt via \Attest, over
  $(c_0,\nu,\pc,\mathit{settled},\mathit{deadline})$;
\item submit the receipt and proof to \Settle;
\item \Settle\ verifies the attestation, public commitments, domain
  anchor and proof, consumes $\nu$, and finalizes the operation.
\end{enumerate}
The construction establishes the cryptographic relation between the
pre-settlement obligation and the post-settlement state.
Section~\ref{sec:security} analyzes the security properties of this
relation.

\section{Security Analysis}
\label{sec:security}

This section characterizes the security guarantees provided by the
settlement-binding construction. The analysis separates three different
claims that are easily conflated:
\begin{itemize}
\item commitment binding, which prevents a commitment from being opened
  to different values;
\item two-phase consistency, which binds the post-settlement state to the
  pre-settlement obligation; and
\item authenticated settlement binding, which connects the resulting
  confidential state to an attestation produced by an authorized PSP.
\end{itemize}
We additionally analyze single-use settlement tags and state the main
settlement-binding theorem.

The analysis uses the following assumptions.

\begin{itemize}
\item \textbf{(A1) Groth16 knowledge soundness.} The deployed Groth16
  system is knowledge-sound for the settlement circuit under its trusted
  setup.
\item \textbf{(A2) Poseidon collision resistance and random-oracle
  privacy heuristic.} It is computationally infeasible to find distinct
  inputs producing the same Poseidon output. For the confidentiality
  discussion only, Poseidon is modeled as a random oracle, with fresh
  high-entropy salts used as blinding values.
\item \textbf{(A3) ECDSA/EIP-712 unforgeability.} Signatures produced
  under the registered PSP keys satisfy EUF-CMA security.
\item \textbf{(A4) Charge-identifier uniqueness.} Within a domain, a
  \txid\ identifies a single payment obligation and is not legitimately
  reused for distinct charges.
\item \textbf{(A5) Keccak-256 collision resistance.} The hash used by
  the attestation oracle to bind the off-chain receipt data is collision
  resistant.
\end{itemize}

We also distinguish these cryptographic and operational assumptions from
the institutional assumption that a registered PSP reports settlement
correctly. The latter cannot be discharged by the cryptographic
construction.

\subsection{Commitment binding}

The first property concerns the individual commitments independently of
the settlement protocol.

\begin{property}[Commitment binding]
No probabilistic polynomial-time adversary can produce two distinct
openings of the same $c_0$ or $c$, except with negligible probability.
\end{property}

For $c_0$, suppose an adversary produces
$(\amount_1,\salt_{0,1},\txid_1)$ and $(\amount_2,\salt_{0,2},\txid_2)$
such that
$\Poseidon(\amount_1,\salt_{0,1},\txid_1) = \Poseidon(\amount_2,\salt_{0,2},\txid_2)$
while the two tuples differ. This constitutes a collision in Poseidon.

The same argument applies to $c$.

Therefore,
\begin{equation}
\Pr[\mathsf{Bad}_{CB}] \le \mathsf{Adv}^{CR}_{\Poseidon}(A). \label{eq:badcb}
\end{equation}
This property is independent of Groth16, the PSP attestation, and the
settlement contract. It is a property of the commitment function itself.

\subsection{Two-phase consistency}

Commitment binding alone is insufficient. The central property of the
construction is that the same confidential obligation is represented
before and after settlement.

Recall that $c_0 = \Poseidon(\amount,\salt_0,\txid)$ and
$c = \Poseidon(\amount,\salt,\eeid,\txid)$.

The circuit contains a single witness variable \amount\ and a single
witness variable \txid, which are used in both equations.

\begin{theorem}[Two-phase consistency]
Assume Groth16 knowledge soundness and Poseidon collision resistance.
Except with negligible probability, every accepting proof contains a
witness for which the \amount\ and \txid\ used in $c$ are identical to
those used in $c_0$.
\end{theorem}

Consequently, if an escrow was created against
$c_0 = \Poseidon(\amount^*,\salt_0^*,\txid^*)$, an accepting proof cannot
use an amount or txid different from $(\amount^*,\txid^*)$ in the
post-settlement commitment, except with negligible probability.

\begin{proof}
Consider an accepting proof $\pi$. By Groth16 knowledge soundness, except
with negligible probability an extractor obtains a witness
$w=(\amount,\salt_0,\salt,\eeid,\txid,k_D,\party,\salt_p)$ satisfying the
circuit relation.

Because the circuit contains one witness variable \amount, that same
value occurs in both commitment equations. Likewise, a single witness
variable \txid\ occurs in both equations. There are therefore no
independent values $\amount_0,\amount_1$ or $\txid_0,\txid_1$ that the
prover can select independently.

It remains to relate the extracted witness to the escrowed commitment. If
$(\amount,\txid) \neq (\amount^*,\txid^*)$, then the extracted witness
and the escrow-opening values define two distinct preimages of the same
$c_0$. This contradicts Poseidon collision resistance.

Hence
\begin{equation}
\Pr[\mathsf{Bad}_{2P}(A)] \le \mathsf{Adv}^{KS}_{\mathrm{Groth16}}(A) +
\mathsf{Adv}^{CR}_{\Poseidon}(A). \label{eq:bad2p}
\end{equation}
\end{proof}

The implementation audit confirms that this equality is enforced
structurally by the circuit: \amount\ and \txid\ are the same private
witness variables in both commitment computations rather than values
compared by an external software check.

This property is deliberately weaker than a statement about real-world
settlement. It establishes consistency between the two confidential
commitments; whether the corresponding external event actually settled
is handled separately by the attestation layer.

\subsection{Authenticated settlement binding}

The previous theorem establishes consistency between the two phases but
does not authenticate the external settlement observation.

The attestation layer adds this missing link.

\begin{property}[Authenticated settlement binding]
An accepted operation is settlement-bound if there exist: a registered
PSP $P$; an unexpired receipt
$R=(c_0,\nu,\pc,\mathit{true},\mathit{deadline})$; a valid PSP signature
$\sigma=\mathsf{Sign}_P(R)$; and a Groth16 witness satisfying the
settlement circuit, such that the receipt and the accepted proof refer to
the same $(c_0,\nu,\pc)$.
\end{property}

This property intentionally does not state that the PSP's assertion
$\mathit{settled}=\mathit{true}$ is objectively true. That fact belongs
to the institutional trust boundary.

\begin{theorem}[Settlement binding]
\label{thm:settlement-binding}
Condition on deterministic correctness of the deployed settlement
contract and immutability of the registered domain anchor $K_D$. Under
assumptions (A1), (A2), (A3), and (A5), for every probabilistic
polynomial-time adversary $A$,
\begin{multline}
\Pr[\mathsf{Bad}_{SB}(A)] \le \mathsf{Adv}^{\mathrm{EUF\text{-}CMA}}_{\mathrm{ECDSA}}(A) \\
+ \mathsf{Adv}^{KS}_{\mathrm{Groth16}}(A) + \mathsf{Adv}^{CR}_{\Poseidon}(A) \\
+ \mathsf{Adv}^{CR}_{\mathrm{Keccak}}(A). \label{eq:badsb}
\end{multline}
\end{theorem}

Here $\mathsf{Bad}_{SB}$ is the event that \Settle\ accepts an operation
for which either: no registered PSP signed the accepted tuple
$(c_0,\nu,\pc,\mathit{true},\mathit{deadline})$, or the accepted tag
$\nu$ is not the value $\Poseidon(\eeid,k_D)$ for the \eeid\ contained in
an accepting witness under the registered domain key.

\begin{proof}[Proof sketch]
Suppose \Settle\ accepts. First, the deployed contract verifies the
Groth16 proof against $(c_0,c,\nu,K_D,\pc)$. By knowledge soundness,
except with negligible probability, an extractor obtains a witness
satisfying the circuit equations. In particular,
$\nu=\Poseidon(\eeid,k_D)$ and $K_D=\Poseidon(k_D)$. If the witness used
a domain key different from the registered key while producing the same
$K_D$, this would yield a Poseidon collision except with negligible
probability.

Second, the contract requires a valid attestation for the same $c_0$,
$\nu$, and \pc. If an adversary produces an accepted receipt without a
registered PSP having signed it, the adversary has forged an
EIP-712/ECDSA signature, contradicting (A3).

Third, the attestation oracle binds its off-chain receipt information to
the corresponding tag using its binding hash. A successful substitution
that preserves the accepted digest while changing the bound values would
contradict the collision resistance of Keccak-256.

Therefore, except with negligible probability, an accepted operation has
both a valid ZK witness and a valid registered-PSP attestation for the
same public tuple.
\end{proof}

The theorem establishes cryptographic consistency and attestation
authenticity. It does not establish that the PSP itself is honest.

\subsection{Single-use settlement}
\label{sec:single-use}

The settlement tag provides replay protection.

For a fixed domain key, $\nu=\Poseidon(\eeid,k_D)$ is deterministic in
the settlement identifier. Consequently, the same \eeid\ produces the
same $\nu$ within that domain.

The public anchor $K_D$ prevents a prover from selecting a different
domain key while still satisfying the same public statement, except with
negligible probability under Poseidon collision resistance.

The reference implementation enforces single use at two independent
points: the settlement contract records accepted tags as consumed
(\code{tagUsed}), and the attestation oracle separately refuses to
attest a tag that it has already attested, regardless of which
commitment the new attestation is for. In the deployment topology used
throughout this paper --- one oracle per deployed settlement instance ---
the oracle's check is the one that actually fires: a second settlement
for the same tag cannot obtain a valid attestation at all, so it never
reaches the tag-consumption check in the settlement contract. The
tag-consumption check in the settlement contract is retained as defense
in depth and becomes load-bearing only in a deployment where several
settlement contracts share a single oracle.

\begin{theorem}[Instance-scoped settlement uniqueness]
\label{thm:instance-scoped}
Fix a domain anchor $K_D$ and a deployed (oracle, \Settle) instance pair.
Under domain-key immutability, Poseidon collision resistance, single-use
enforcement of $\nu$, and the charge-identifier uniqueness assumption, no
probabilistic polynomial-time adversary can cause two accepted
settlements corresponding to the same \eeid, except with negligible
probability.
\end{theorem}

\begin{proof}[Proof sketch]
Suppose two settlements correspond to the same \eeid. Because the domain
key is fixed, $\Poseidon(\eeid,k_D)$ produces the same $\nu$ for both.
The public anchor $K_D$ prevents the two accepting proofs from using
unrelated domain keys except with negligible probability. The first
accepted settlement is attested and consumes the tag at the oracle. The
second settlement, for a different commitment but the same tag, cannot
obtain an attestation from that oracle; absent an attestation it cannot
satisfy the finalize precondition that requires one. Hence two accepted
settlements for the same \eeid\ cannot occur within the same deployed
(oracle, \Settle) instance pair, except with negligible probability.
\end{proof}

This was confirmed experimentally in the reference implementation: a
second settlement attempt over a different commitment carrying the same
tag is rejected by the oracle's attestation call itself, before the
settlement contract's own finalize check is reached.

The scope qualifier is important. The reference implementation does not
provide a globally shared spent-tag registry, at either the oracle or the
settlement-contract level. Consequently, the property does not extend
automatically across two independently deployed (oracle, \Settle)
instance pairs, nor across domains with different keys. This was also
confirmed experimentally: two independently deployed (oracle, \Settle)
pairs, each given the same commitment, tag, and proof, both accept the
settlement, because neither oracle has prior knowledge of the other's
attestation. A shared oracle, or a shared registry indexed by $K_D$,
would be required to elevate this property from instance scope to domain
scope.

\subsection{Attestation authenticity versus settlement truth}

The construction deliberately separates two statements that are often
conflated: ``a registered PSP signed `settled'\," and ``the external
payment rail actually settled.'' The first is a cryptographic property.
The second is an institutional property.

\begin{property}[Attestation authenticity]
If \Settle\ accepts an operation, then, except with negligible
probability under ECDSA/EIP-712 EUF-CMA security, a registered attestor
produced the accepted receipt.
\end{property}

This follows directly from signature verification, registration
checking, domain separation, and expiry checking.

However, a malicious or compromised registered PSP can still produce a
valid signature for a false settlement claim. No cryptographic reduction
in this construction eliminates that possibility.

We therefore model the remaining trust as:
\begin{quote}
\small
\textbf{(I1)} Registered PSP honestly reports observed settlement.
\end{quote}
Under (I1), the cryptographic settlement-binding theorem can be
interpreted operationally as:
\[
\begin{aligned}
\mathit{Accepted} \;\Rightarrow\;\; &\mathit{ConfidentialStateConsistent} \\
\wedge\;\; &\mathit{PSPAuthenticated} \\
\wedge\;\; &\mathit{SettlementObserved}_{PSP}.
\end{aligned}
\]
Without (I1), only the first two terms are established.

This distinction is essential to the scope of the construction and
prevents the protocol from being interpreted as a zero-knowledge proof of
central-bank or payment-rail truth. The original analysis explicitly
identifies this residual PSP assumption as outside the cryptographic
reduction.

\subsection{Main settlement-binding theorem}

We can now state the principal security claim of the construction.

\begin{theorem}[Cryptographic settlement binding]
Assume: Groth16 knowledge soundness for the deployed settlement circuit;
Poseidon collision resistance; ECDSA/EIP-712 EUF-CMA security; Keccak-256
collision resistance for the attestation binding hash; uniqueness of the
payment-obligation identifier (\txid); an immutable domain anchor $K_D$;
and deterministic correctness of the deployed \Settle\ contract. Then
every accepted settlement is, except with negligible probability,
cryptographically consistent with the pre-settlement commitment and
authenticated as the same $(c_0,\nu,\pc)$ tuple by a registered PSP. More
precisely,
\begin{multline}
\Pr[\mathsf{Bad}_{SB}(A)] \le \mathsf{Adv}^{\mathrm{EUF\text{-}CMA}}_{\mathrm{ECDSA}}(A) \\
+ \mathsf{Adv}^{KS}_{\mathrm{Groth16}}(A) + \mathsf{Adv}^{CR}_{\Poseidon}(A) \\
+ \mathsf{Adv}^{CR}_{\mathrm{Keccak}}(A). \label{eq:mainbound}
\end{multline}
\end{theorem}

The operational assumption on \txid\ uniqueness is required to interpret
the identifier as a unique payment obligation, while the institutional
assumption that the PSP truthfully reports settlement remains outside the
cryptographic bound.

The theorem therefore establishes the following implication:
\begin{quote}
\small
$\mathsf{SettleAccept} \Rightarrow \exists w$ such that
$c_0=\Commit_0(w)$, $c=\Commit(w)$, $\nu=\mathsf{Tag}(w)$, and
$\pc=\mathsf{PartyCommit}(w)$; and there exists a registered PSP $P$ such
that $P$ signed $(c_0,\nu,\pc,\mathit{true},\mathit{deadline})$.
\end{quote}
The theorem does not claim that the payment amount, payment identifier,
or counterparty is publicly recoverable from the accepted transcript. Nor
does it claim economic irreversibility of the underlying payment. Those
are separate properties considered in the privacy and limitations
analyses.

The significance of the construction is therefore not that it removes
the PSP from the trust model. Rather, it converts an external settlement
observation into a cryptographically constrained statement about a
confidential on-chain state transition, while preserving the
confidentiality of the underlying settlement data.

\section{Privacy Analysis}
\label{sec:privacy}

The settlement-binding construction is designed to hide the sensitive
attributes of the external payment while preserving publicly verifiable
evidence that the settlement is bound to the confidential state
transition.

The privacy claim is deliberately narrower than anonymity or
unlinkability.

\subsection{Confidential values}

The following values are private witness elements: \amount, \eeid,
\txid, \party. They are never supplied as public inputs to the
settlement verifier.

The public transcript instead contains $c_0,c,\nu,K_D,\pc$, together with
the zero-knowledge proof and the PSP's signed receipt.

The amount appears only through the commitments $c_0$ and $c$, where the
salts provide high-entropy blinding.

Likewise, the counterparty appears only through
$\pc=\Poseidon(\party,\salt_p)$, and the payment identifier appears only
through the keyed tag $\nu=\Poseidon(\eeid,k_D)$.

For this confidentiality discussion, we use a random-oracle heuristic in
which the fresh high-entropy salts act as blinding values. Together with
the zero-knowledge property of Groth16, this prevents the public
transcript from directly revealing the corresponding witness values. This
is a privacy heuristic rather than a standalone composable hiding theorem
for Poseidon.

The PSP receipt does not contain the amount, \eeid, \txid, or plaintext
counterparty identifier. It therefore does not introduce a plaintext
disclosure channel.

\subsection{What the construction does not hide}

The construction does not provide complete transaction anonymity. In
particular, the following remain public or potentially observable:
existence of the on-chain operation; timing of the settlement
transaction; submitting and relaying addresses; the registered attestor
address; the domain anchor $K_D$; repeated counterparty commitments
(\pc); ordinary transaction-level metadata.

Therefore, confidentiality $\neq$ unlinkability.

An observer may be unable to recover the amount or payment identifier
while still correlating multiple operations through public metadata.

This is intentional. The construction provides confidential settlement
binding, not an anonymity system.

\subsection{Selective linkability of the settlement tag}

The keyed tag $\nu=\Poseidon(\eeid,k_D)$ provides a more specific privacy
property.

An observer outside the domain who does not know $k_D$ cannot directly
compute candidate tags for known \eeid's. In contrast, an institution
holding $k_D$ can compute $\Poseidon(\eeid,k_D)$ for \eeid\ values it
legitimately observes and therefore correlate those payments with
on-chain tags.

The tag should therefore be understood as a selective-linkability
mechanism, rather than as a universal privacy mechanism.

This distinction is important because the PSP and proving institutions
are already trusted participants in the payment domain and necessarily
possess information that an external blockchain observer does not.

\subsection{Counterparty commitments}

The counterparty commitment provides a reusable pseudonymous anchor for
the settlement participant: $\pc=\Poseidon(\party,\salt_p)$.

Its purpose is not to prove the legal identity of a person or
organization on-chain. Rather, it provides a confidential reference
against which zero-knowledge policy predicates may subsequently be
evaluated.

The commitment is meaningful only because the PSP computes it from the
debtor identifier contained in the payment message and the finalization
procedure requires the attested \pc\ to equal the public \pc\ of the
settlement proof.

Without this cross-check, a prover could select a favorable counterparty
independently of the actual payer.

\subsection{Privacy boundary}

\begin{table*}[t]
\centering
\begin{tabular}{@{}ll@{}}
\toprule
Property & Status \\
\midrule
Amount confidentiality & Protected \\
\eeid\ confidentiality & Protected \\
\txid\ confidentiality & Protected \\
Counterparty identifier confidentiality & Protected \\
Operation existence & Not protected \\
Transaction timing & Not protected \\
Submitter/relayer privacy & Not protected \\
PSP identity privacy & Not protected \\
Cross-operation unlinkability & Not guaranteed \\
\bottomrule
\end{tabular}
\caption{Privacy boundary of the settlement-binding construction.}
\end{table*}

The construction consequently claims confidentiality of settlement
attributes, not anonymity of the settlement transaction itself.

\section{Pix Instantiation}
\label{sec:pix}

The general settlement-binding construction is instantiated using Pix as
the external settlement rail.

The abstraction requires three rail-side capabilities: a
payment-obligation identifier available before settlement; a
settlement-specific identifier available to an authorized observer of
the payment; and an attestor capable of confirming the settlement and
recomputing the relevant commitment.

For Pix \cite{spicatalog}, the payment-obligation identifier is the
merchant-side \txid, while the settlement-specific identifier is the
EndToEndId (\eeid). A regulated PSP acts as the attestor.

The distinction between the two identifiers is important.

The \txid\ belongs to the charge and can therefore participate in the
pre-settlement commitment: $c_0=\Poseidon(\amount,\salt_0,\txid)$.

The \eeid\ is associated with the actual payment transaction and
therefore enters the post-settlement commitment:
$c=\Poseidon(\amount,\salt,\eeid,\txid)$.

The construction does not rely on the identifier itself being evidence
of settlement. An \eeid\ can exist for a payment that is subsequently
rejected. Settlement status is therefore supplied by the PSP attestation
rather than inferred from possession of the identifier.

This makes the Pix instantiation representative of a broader class of
instant-payment systems. The same architecture could apply to another
rail if it exposes a unique settlement identifier, a pre-settlement
obligation identifier, and an authorized institution capable of
attesting the resulting settlement.

\subsection{Settlement semantics}

The reference implementation defines $\mathit{settled}=\mathit{true}$ to
mean that the receiver has credit availability, rather than merely that
the payment has cleared an intermediate rail state.

This distinction matters because a payment can remain subject to return
or fraud-recovery mechanisms after credit becomes available.

The protocol therefore distinguishes: \Finalized\ (the cryptographic
settlement-binding conditions were satisfied) and \Irrevocable\ (the
protocol's configured dispute window has elapsed without an accepted
reversal).

\Irrevocable\ is a protocol state and must not be interpreted as a claim
of economic irreversibility by the payment rail.

The reference deployment uses a 72-hour dispute window. The underlying
rail may permit certain return or fraud-recovery mechanisms for
substantially longer periods. Consequently, an application requiring
economic finality must choose its own risk policy rather than treating
the protocol state as a substitute for the rail's legal or operational
finality.

\section{Policy Predicates as an Application Layer}
\label{sec:predicates}

The settlement-binding primitive can be consumed by policy predicates
without exposing the underlying settlement data.

These predicates are intentionally treated as an application layer, not
as part of the core settlement-binding theorem.

Three reference predicates are considered.

\subsection{Value bound}

A range circuit proves that the committed settlement amount is below a
policy limit: $\amount < \mathit{limit}$.

The predicate uses a bounded comparison over the opening of the
settlement commitment.

The statement is deliberately called a policy predicate rather than a
universal compliance determination. It establishes only that the amount
satisfies the specified limit under the committed value.

\subsection{Sanctions-set non-membership}

A Sparse Merkle Tree represents an authority-defined set of sanctioned
identifiers.

The prover demonstrates that the committed counterparty does not occur
in the tree: $\party \notin S$.

The public statement consists of the counterparty commitment and the
policy root.

This does not establish that the counterparty is legally ``unsanctioned''
in an absolute sense. The result is relative to the particular root, its
version, its normalization procedure, and the authority responsible for
maintaining it.

\subsection{KYC-set membership}

The dual predicate proves $\party \in \mathsf{KYCRoot}$.

Again, the cryptographic statement is set membership. It does not
independently establish the correctness, freshness, jurisdictional
scope, or lifecycle of the underlying KYC process.

\subsection{Binding predicates to the actual payer}

The important architectural point is that the predicates operate on the
settlement's attested \pc, rather than on an arbitrary prover-selected
identity.

The settlement protocol first establishes $\pc=\Poseidon(\party,\salt_p)$
and requires the PSP attestation to contain the same \pc.

A consumer of a policy predicate must then verify that its public
commitment corresponds to the \pc\ of the settled operation.

Thus the composition is: PSP-attested payer $\rightarrow$ \pc\
$\rightarrow$ ZK policy predicate.

Without this composition check, a prover could demonstrate a predicate
about a different identity.

The predicates are therefore useful as an extension of the
settlement-binding primitive, but they do not enlarge the cryptographic
claim of the core construction.

\section{Implementation and Evaluation}
\label{sec:impl}

\subsection{Reference implementation}

The settlement circuit is implemented in Circom 2.2.1, the contracts use
Solidity 0.8.27, and the Groth16 verifier is generated with snarkjs
0.7.5, matching the toolchain versions pinned in the reference
repository.

The primary circuit is \code{settlement\_verify\_v2}.

At optimization level \code{-O2}, it contains 1,251 R1CS constraints. At
\code{-O1}, the circuit contains 1,266 nonlinear and 1,524 linear
constraints.

The circuit exposes five public signals and eight private witness
values.

The circuit contains five Poseidon gadgets with arities 3, 4, 2, 1, 2,
corresponding to the two commitments, settlement tag, domain anchor, and
counterparty commitment.

\subsection{On-chain cost}

The primary figure is the complete settlement transaction rather than an
isolated verifier call.

On the XDC Apothem public testnet, chain ID 51, the measured transaction
costs were:

\begin{table*}[t]
\centering
\begin{tabular}{@{}lr@{}}
\toprule
Operation & Gas \\
\midrule
\code{lock(c0, deadline)} & 74{,}982 \\
\code{attest(receipt, signature)} & 90{,}913 \\
\code{finalize(proof)} & 353{,}062 \\
isolated \code{verifyProof} & 243{,}182 \\
\bottomrule
\end{tabular}
\caption{Measured XDC Apothem gas costs.}
\end{table*}

These four figures are reconfirmed against the recorded Apothem
deployment log and the automated test suite. \code{cancel(c0)} and
\code{attestReversal}, which appeared in an earlier draft of this table,
are omitted here because they were not reconfirmed against the Apothem
deployment; a local-network reproduction gives 28{,}115 and 57{,}541 gas
respectively, close to but not identical to the previously recorded
Apothem figures (30{,}259 and 63{,}529), and the discrepancy has not been
attributed to a specific cause.

The complete \code{finalize} transaction is the primary deployment
metric because it includes proof verification, attestation
cross-checking, domain-anchor validation, tag consumption, and state
transition.

The isolated verifier transaction includes its own base transaction cost
and proof calldata and therefore cannot simply be subtracted from the
complete \code{finalize} transaction to derive a verifier-only execution
cost.

Groth16 pairing verification remains the dominant component of the
settlement cost.

\subsection{Circuit well-determinedness}

The circuits were additionally analyzed with Ecne \cite{ecne} for
well-determinedness.

This analysis is explicitly not presented as formal verification of
functional correctness. It addresses the narrower question of whether
circuit signals are sufficiently constrained.

The meaningful \code{-O0} results were:

\begin{table*}[t]
\centering
\begin{tabular}{@{}lccc@{}}
\toprule
Circuit & Targets solved & Signals & Result \\
\midrule
\code{settlement\_verify\_v2\_out} & 5/5 & 4{,}229/4{,}229 & Sound \\
\code{kyc\_inclusion\_verify\_out} & 1/1 & 18{,}230/18{,}248 & Sound \\
\bottomrule
\end{tabular}
\caption{Ecne \code{-O0} well-determinedness results.}
\end{table*}

For the optimized settlement circuit, Ecne solved 4/5 declared targets
and therefore returned inconclusive, not unsound.

This distinction is important: the optimized result does not constitute
a negative security result. It indicates that the symbolic analysis
could not determine one target under the optimized constraint
representation.

The analysis also exposed a methodological pitfall in which circuits
whose public outputs are expressed only as input-plus-equality
constraints can produce a vacuous Ecne verdict. The paper therefore
reports the number of targets explicitly rather than relying on the
tool's raw ``sound'' output.

The well-determinedness analysis establishes the absence of a particular
class of under-constrained-circuit bugs. It does not establish
functional conformance with the intended protocol specification.

\subsection{End-to-end deployment}

The construction was deployed and exercised on the XDC Apothem public
testnet.

The complete flow was executed: lock $c_0$ $\rightarrow$ PSP receipt
$\rightarrow$ attestation $\rightarrow$ Groth16 proof $\rightarrow$
finalize.

No amount, \eeid, or charge identifier was included in the corresponding
settlement calldata.

Negative on-chain tests exercised: non-attestor signatures,
$\mathit{settled}=\mathit{false}$, expired receipts, mismatched $c_0$,
mismatched counterparty commitment, incorrect domain key, reversal, and
premature and post-deadline cancellation. These cases are exercised by a
manually run script against the Apothem deployment and are not part of
the automated test suite.

Each negative case reverted as intended in the reference deployment.

Tag replay is additionally covered by an automated test against a local
deployment of the same contracts, using two independently generated
Groth16 proofs that share an \eeid\ and domain key and therefore share a
settlement tag while committing to different terms. A second settlement
for the same tag is rejected, and --- as discussed in
Section~\ref{sec:single-use} --- the rejection occurs at the attestation
oracle rather than at the settlement contract's own tag-consumption
check, which only becomes load-bearing when several settlement contracts
share one oracle. The same automated test also confirms, by constructing
two independently deployed (oracle, \Settle) pairs, that this single-use
property does not extend across independent deployments, consistent with
the instance-scoped qualifier in Theorem~\ref{thm:instance-scoped}.

The reproducibility claims in this section refer to the repository state
identified by commit \code{0db97bf91ddc5be475e3d4175a62e241eb80df79}.

The prototype deployment uses the same key for deployment, attestor
registration, and registry administration. This is explicitly an
implementation convenience and not a design requirement. A production
deployment should separate these roles and use appropriate governance
controls.

\section{Related Work}
\label{sec:related}

\subsection{External payment evidence}

The closest prior work is ZKP2P \cite{zkp2p}, which demonstrates that
external payment evidence can be used to authorize an on-chain action.

The present construction does not claim novelty for the general pattern
``external payment evidence $\rightarrow$ on-chain action.''

Its target is narrower: pre-settlement confidential commitment
$\rightarrow$ external settlement observation $\rightarrow$ confidential
on-chain state.

The distinguishing feature is therefore the temporal and cryptographic
relationship between the state committed before settlement and the
external settlement identifier available afterward.

The attestation model also differs: the reference construction uses a
regulated PSP as the institutional observer rather than consumer-facing
email or TLS evidence.

\subsection{Confidential EVM transfers}

Systems such as Zeto \cite{zeto}, Nightfall, Railgun, and Rayls/Enygma
\cite{rayls,rayls2} address confidential transfer validity.

Their core question is approximately: ``Is this private on-chain state
transition internally valid?''

The settlement-binding construction addresses a different question:
``Does a specific externally observed settlement correspond to this
confidential on-chain state?''

The construction is therefore complementary to confidential asset
systems rather than a replacement for them.

In particular, the resulting commitment can serve as an upstream input
to a confidential asset protocol.

No direct integration with Zeto or Paladin is claimed in this paper.

\subsection{Institutional privacy and delivery-versus-payment}

Canton \cite{canton} and the Drex ecosystem address institutional
privacy and delivery-versus-payment in permissioned or regulated
infrastructures.

The reference construction instead places the settlement-binding layer
on a public EVM chain while using an institutional PSP as the attestor.

The resulting trust boundary is therefore different: public EVM +
regulated attestation + zero-knowledge confidentiality.

\subsection{Confidential compliance and auditability}

Zeto \cite{zeto} and Qurrency \cite{qurrency} provide related mechanisms
for confidential compliance and auditor access.

The policy predicates in this work differ primarily in their disclosure
semantics: the consumer can learn a boolean predicate result without
necessarily receiving the underlying amount or identity.

The post-quantum audit channel is similarly complementary rather than
central to the settlement-binding contribution.

\subsection{Positioning}

The scientific contribution should therefore not be stated as a
collection of novel cryptographic primitives.

Poseidon commitments \cite{poseidon}, Groth16 \cite{groth16}, EIP-712
signatures, nullifier-like tags, Sparse Merkle Trees, and confidential
transfer mechanisms are individually established.

The contribution is their composition around a specific missing
interface: external settlement event $\rightarrow$ cryptographic
settlement binding $\rightarrow$ confidential EVM state.

This positioning is also why the Pix implementation is an instantiation
of the construction rather than the construction itself.

\section{Limitations and Open Problems}
\label{sec:limitations}

\subsection{Trusted PSP}

The largest remaining trust assumption is the PSP.

The construction proves that a registered PSP authenticated the accepted
statement. It does not prove that the PSP reported truthfully.

The strongest future direction is therefore to replace or weaken this
assumption through one of several approaches: a central-bank or
rail-level signature; signed rail artifacts verified retrospectively;
dual attestation from both sides of the payment; zkTLS/DECO
\cite{deco} over an authoritative PSP endpoint; or a periodically signed
accumulator of settled payment identifiers.

Of these, retrospective rail anchoring is particularly attractive
because it does not have to remain on the critical path of settlement.

A signed membership root could allow an institution to prove, after the
fact, that its settlement identifier belongs to the rail's settled set
while still exposing only the keyed tag on-chain.

Such a mechanism, however, requires an explicit completeness guarantee
for the published root. Without completeness, non-membership in a
published root does not prove non-settlement.

\subsection{Instance-scoped uniqueness}

The current spent-tag registry is scoped to the deployed (oracle,
\Settle) instance.

A production domain spanning multiple settlement contracts would benefit
from a shared registry indexed by $K_D$.

This would elevate the uniqueness guarantee from instance scope to
domain scope.

\subsection{Domain-key compromise and rotation}

The domain key $k_D$ is a security boundary.

A compromised key can reveal the correspondence between observed \eeid\
values and historical settlement tags. It does not, by itself, forge PSP
attestations or Groth16 proofs.

Key rotation is also non-trivial. Rotating $k_D$ produces a new tag for
the same \eeid\ and therefore cannot preserve the current spent-tag
uniqueness property across domains automatically.

A production system must consequently treat key rotation as a domain
transition rather than a transparent key-management operation.

\subsection{Trusted setup}

The reference settlement circuit uses a Groth16 setup with a single
contributor followed by a beacon.

This is sufficient for the experimental artifact but is not a
production-grade ceremony.

A production deployment should use a multi-party ceremony with
independent contributors and a verifiable public randomness beacon.

The trusted setup limitation affects the soundness guarantees relying on
Groth16; it does not affect the independent commitment-binding or
signature-unforgeability properties.

\subsection{Economic finality}

The protocol's \Irrevocable\ state is not equivalent to economic
finality on the underlying payment rail.

The reference deployment uses a 72-hour dispute window, while certain
rail-level return and fraud mechanisms can remain effective for
substantially longer periods.

A deployment requiring stronger asset protection must either extend its
risk window or introduce a rail-level anchoring mechanism that can
represent later reversals.

\subsection{Atomic delivery-versus-payment}

The two-phase construction provides an escrow commitment before
settlement, which is the necessary ingredient for DvP composition.

However, a complete atomic composition with a confidential asset system
has not yet been implemented and measured.

An integration with a system such as Zeto or Paladin would therefore be
a meaningful next experiment rather than a result claimed by this paper.

\subsection{Benchmark methodology}

This paper reports on-chain gas costs but does not report off-chain
proving and verification time. An earlier internal measurement pass was
excluded because it consisted of single warm runs rather than a
statistically meaningful benchmark.

A proper benchmark should report proving and verification time over
$N \geq 30$ warm runs per circuit, with median and dispersion, on fixed
and disclosed hardware. This is left for future work.

Likewise, verifier comparisons across circuits should use the same
generated verifier template before attributing gas differences to the
number of public signals.

\subsection{Well-determinedness analysis}

The Ecne analysis provides evidence against one class of
under-constrained circuits but does not prove functional correctness.

Future work should cross-check the result with an independent tool such
as $\mathrm{QED}^2$ \cite{qed2}, identify the currently unresolved
signals in the KYC circuit, and extend the analysis to the optimized
circuit representation.

\subsection{Post-quantum audit disclosure}

The ML-KEM \cite{mlkem} audit channel remains preliminary.

The full in-circuit ML-KEM binding compiles at approximately 1.62
million constraints, but end-to-end proving and on-chain verification
measurements have not yet been established.

It should therefore remain outside the core security claim until those
measurements and artifacts are complete. ML-KEM is not part of the
security argument for cryptographic settlement binding; it is explored
here solely as a preliminary post-quantum audit-disclosure extension.

\section{Conclusion}

We introduced a cryptographic settlement-binding construction for
confidential EVM state.

The central problem is not to prove that a payment identifier exists,
nor to replace an institutional settlement observer with zero knowledge.
The problem is to establish a cryptographic correspondence between two
events that occur at different times: a confidential obligation
committed on-chain before external settlement; and an external
settlement event observed and authenticated by an authorized
institution.

The construction addresses this problem through a two-phase commitment.

Before settlement,
\[
c_0=\Poseidon(\amount,\salt_0,\txid)
\]
provides an on-chain commitment to the agreed terms.

After settlement,
\[
c=\Poseidon(\amount,\salt,\eeid,\txid)
\]
incorporates the settlement-specific identifier.

A Groth16 relation proves that both commitments share the same
confidential amount and charge identifier. A domain-keyed tag
$\nu=\Poseidon(\eeid,k_D)$ provides deterministic, instance-scoped
single-use identification without exposing the underlying payment
identifier. A registered PSP then authenticates the tuple
$(c_0,\nu,\pc,\mathit{true},\mathit{deadline})$ through an EIP-712
receipt.

The resulting security statement is deliberately precise:
\begin{quote}
\small
\emph{accepted settlement} $\Rightarrow$ \emph{cryptographically
consistent confidential state} $+$ \emph{authenticated PSP statement.}
\end{quote}
The construction does \textbf{not} claim to establish central-bank or
payment-rail truth without trusting the attesting institution. Removing
or weakening this assumption remains the principal open problem.

The Pix instantiation demonstrates that the construction can be
implemented on a public EVM. The primary settlement circuit contains
1,251 R1CS constraints at \code{-O2}, while the reference \code{finalize}
transaction costs 353,062 gas on the XDC Apothem testnet. The
implementation also demonstrates negative acceptance cases and provides
circuit well-determinedness evidence under the stated methodology.

More broadly, the construction is intended as an interface between two
otherwise separate systems:
\begin{quote}
\small
\emph{real-world settlement} $\rightarrow$ \emph{settlement binding}
$\rightarrow$ \emph{confidential asset state.}
\end{quote}
This interface can be instantiated for Pix and potentially for other
instant-payment rails, and can serve as an upstream component for
confidential asset systems.

The principal result is therefore not a new commitment primitive or a
new zero-knowledge system. It is a formally characterized method for
cryptographically binding an externally observed settlement event to a
confidential on-chain state transition while revealing only
application-required information.

\appendix

\section{Pix Canonicalization Profile}
\label{app:pix}

This appendix specifies the concrete canonicalization profile used by
the Pix instantiation of Section~\ref{sec:pix}. It is a deployment
profile and is not part of the cryptographic settlement-binding
construction itself.

\subsection{Counterparty Identifier}

For CPF, the canonical representation is the 11 decimal digits of the
identifier, with punctuation removed and no padding.

For CNPJ, the canonical representation is the 14-character identifier,
upper-cased and with punctuation removed.

The resulting representation is encoded as the exact byte sequence
consumed by the commitment function.

\subsection{\eeid\ and \txid}

For the Pix deployment, \eeid\ and \txid\ are represented by the exact
bytes returned by the PSP API.

No case folding, trimming, or re-serialization is performed.

This rule avoids ambiguity in the byte representation used by the
commitment function.

\subsection{Provenance}

The \eeid\ is generated by the payer's PSP when the payment is
initiated. It can therefore exist for a transaction that is later
rejected as well as for one that settles.

Consequently, the presence of an \eeid\ is not treated as evidence of
settlement. The settlement status is supplied through the authorized PSP
attestation described in Sections~\ref{sec:model}--\ref{sec:security}.

\subsection{Scope}

Nothing in the core construction depends on these Pix-specific
representations. Another instant-payment rail can replace the
canonicalization profile with its own rail-defined identifier
representation while preserving the settlement-binding construction.

\section{Threat Model Matrix}
\label{app:threats}

\begin{table*}[t]
\centering
\scriptsize
\setlength{\tabcolsep}{2pt}
\begin{tabular}{@{}p{3.05cm}p{2.15cm}p{4.5cm}p{4.5cm}@{}}
\toprule
Adversary / Failure & Protected Property & Mitigation & Residual Limitation \\
\midrule
False proof & ZK consistency & Groth16 knowledge soundness & Depends on the trusted setup \\
Different amount in $c_0$ and $c$ & Two-phase binding & Same private amount is constrained in both commitments & None under the stated circuit and hash assumptions \\
Different operation identifier in $c_0$ and $c$ & Operation binding & Same private \txid\ is constrained in both commitments & Assumes identifiers are not legitimately reused across charges \\
Reuse of a settlement tag & Single use & Write-once oracle attestation, with \code{tagUsed} as contract-level defense in depth & Scoped to the deployed (oracle, \Settle) instance pair \\
Replay across domains & Domain separation & $K_D=\Poseidon(k_D)$ and $\nu=\Poseidon(\eeid,k_D)$ & Global uniqueness across domains is not provided \\
Forged PSP receipt & Attestation authentication & EIP-712 + ECDSA recovery + registered-attestor check & Does not protect against a malicious registered PSP \\
Expired receipt & Temporal validity & Receipt deadline checked on-chain & Deadline policy is deployment-specific \\
$\mathit{settled}=\mathit{false}$ receipt & Settlement acceptance & Finalization requires $\mathit{settled}=\mathit{true}$ & Truthfulness remains an institutional assumption \\
Wrong $c_0$ in receipt & Value/operation consistency & Receipt $c_0$ must equal locked $c_0$ & A malicious PSP can still sign an incorrect statement \\
Wrong counterparty commitment & Counterparty binding & Receipt \pc\ must equal proof \pc\ & Correctness depends on the PSP's attestation \\
Prover chooses a favorable counterparty & Predicate binding & \pc\ is a public circuit signal and signed receipt field & Consumer must check the predicate against the settled operation \\
\eeid\ alone presented as settlement evidence & External-event authentication & PSP attestation required & No direct rail proof exists in the construction \\
Chain observer learns amount & Confidentiality & Amount appears only inside ZK witness and commitments & Side channels and external information may still correlate operations \\
Chain observer learns \eeid & Confidentiality & \eeid\ appears only inside $c$ and $\nu$ & Domain-key holders can deliberately link tags to \eeid\ values \\
Domain key compromise & Selective linkability & Key is never public & Historical tags in that domain may become linkable \\
Per-operation salt compromise & Commitment hiding & Fresh salts & Exposure affects the corresponding operation \\
PSP signing-key compromise & Attestation integrity & Key revocation and receipt deadlines & Attacker may forge attestations during the usable window \\
PSP refuses to attest & Liveness & Lock timeout and cancellation & Settlement cannot be finalized without attestation \\
Front-running \code{finalize} & Operation ownership & \code{finalize} changes the already-selected operation and does not assign benefit to \code{msg.sender} & Downstream applications must preserve this invariant \\
External refund after finalization & Economic finality & Configurable reversal/dispute window & Rail-level return mechanisms may outlive the protocol window \\
Registry-owner compromise & Governance & Administrative controls & Production deployment should use stronger governance \\
Trusted-setup compromise & Proof soundness & Production ceremony required & Reference deployment uses a research-prototype ceremony \\
Central-bank/rail false statement & External-world truth & Not cryptographically mitigated & Requires stronger rail-level evidence or additional attestors \\
\bottomrule
\end{tabular}
\caption{Threat model matrix.}
\end{table*}

\subsection{Trust Decomposition}

The construction intentionally decomposes the final claim into
independently attributable components.

\begin{table*}[t]
\centering
\scriptsize
\setlength{\tabcolsep}{3pt}
\begin{tabular}{@{}p{7.5cm}p{7.5cm}@{}}
\toprule
Claim & Evidence \\
\midrule
Same amount and operation identifier occur in both commitments & Zero-knowledge proof \\
Settlement tag is derived from the settlement identifier & Zero-knowledge proof \\
Settlement tag is single-use & Oracle write-once attestation + on-chain \code{tagUsed} defense in depth \\
\eeid\ is bound to the accepted confidential state & Zero-knowledge proof + PSP attestation \\
Counterparty commitment is the one accepted for settlement & PSP attestation + on-chain cross-check \\
External settlement actually occurred & PSP attestation under institutional assumption $I_1$ \\
Central-bank / rail infrastructure itself attests the event & \textbf{Not established by this construction} \\
Operation cannot remain locked forever & Protocol timeout \\
Settlement is economically irreversible & \textbf{Not established cryptographically} \\
\bottomrule
\end{tabular}
\caption{Trust decomposition.}
\end{table*}

This decomposition is the security boundary used by the formal analysis
in Section~\ref{sec:security}.

\section{Ecne Constraint-Analysis Methodology}
\label{app:ecne}

This appendix documents the methodology used to analyze the
zero-knowledge circuits with Ecne. The purpose is to make the reported
results reproducible and to prevent an overly strong interpretation of
the tool's verdict.

\subsection{Scope of the Analysis}

Ecne is used here to analyze whether the generated R1CS is
well-determined, i.e., whether the signals relevant to the declared
targets are uniquely constrained by the system.

This is an under-constraint analysis.

It is \textbf{not} a proof of:
\begin{itemize}
\item functional correctness of the circuit;
\item correctness of the protocol specification;
\item cryptographic security of Poseidon or Groth16;
\item correctness of the Solidity verifier;
\item truthfulness of the PSP attestation; or
\item security of the complete application.
\end{itemize}
Accordingly, this paper refers to the result as \textbf{constraint
analysis}, rather than formal verification.

This analysis is also distinct from static-analysis linting of the
Circom source itself, such as Circomspect \cite{circomspect}, which
flags coding-pattern risks (for example, unconstrained signals
introduced by common template misuse) without solving the resulting
R1CS. The two approaches are complementary: a circuit can pass a linter
and still be under-constrained, or fail a linter check while remaining
fully constrained.

\subsection{Avoiding the Vacuous-Target Problem}

A direct Ecne run over a circuit written with public inputs constrained
through equality statements can produce a misleadingly positive result.

In the original circuit structure, a public value may be supplied as an
input and constrained with a relation such as
\[
c \mathrel{===} H(w).
\]
Depending on compilation, the resulting R1CS can expose no public
outputs to the Ecne target set. If the target set is empty, the
criterion
\[
\mathit{targets\ solved} = \mathit{targets}
\]
reduces to $0=0$. The resulting ``sound'' verdict is therefore vacuous.

We therefore do not report the raw verdict alone.

\subsection{Corrected Analysis}

For the analysis reported in the paper, dedicated \code{\_out} variants
of the circuits were used.

The derived public signals were declared as explicit outputs, while
preserving the constraint semantics of the original circuits. The
analysis was performed on the unoptimized \code{-O0} R1CS because the
optimizer can fold linear structure that Ecne's symbolic analysis relies
upon.

The primary settlement circuit was analyzed as
\code{settlement\_verify\_v2\_out}, and the KYC inclusion circuit as
\code{kyc\_inclusion\_verify\_out}.

\subsection{Results}

\begin{table*}[t]
\centering
\small
\begin{tabular}{@{}lcccl@{}}
\toprule
Circuit & Opt. & Declared & Solved & Verdict \\
\midrule
\code{settlement\_verify\_v2\_out} & O0 & 5 & 5 & sound \\
\code{kyc\_inclusion\_verify\_out} & O0 & 1 & 1 & sound \\
\code{settlement\_verify\_v2\_out} & O2 & 5 & 4 & inconclusive \\
\bottomrule
\end{tabular}
\caption{Ecne results by circuit and optimization level. Signals:
4{,}229/4{,}229 (O0 settlement), 18{,}230/18{,}248 (O0 KYC),
1{,}016/1{,}260 (O2 settlement).}
\end{table*}

The O0 settlement result solves all five declared targets and all
reported signals.

The KYC circuit has one declared target, which is solved. The remaining
18 signals reported by Ecne as undetermined are therefore not among the
declared targets. Those signals have not yet been individually mapped
against the generated symbol file.

The O2 settlement analysis is reported as \textbf{inconclusive}, not
unsound. Ecne solved four of the five targets but could not determine
the remaining target under the optimized constraint system.

\subsection{Interpretation}

The reported results provide evidence against one important class of
implementation error: under-constrained R1CS signals.

They do not establish that the circuit computes the intended predicate.

For example, a circuit can be fully determined while computing an
incorrect hash relation. Functional correctness therefore remains
established by the circuit specification, manual audit, implementation
cross-checks, and the cryptographic arguments of
Section~\ref{sec:security} rather than by Ecne alone.

A different under-constraint analysis, such as the techniques
represented by $\mathrm{QED}^2$ \cite{qed2}, would provide an independent
cross-check. Such an analysis has not been included in the present
evaluation.

\subsection{Relationship to the Main Security Argument}

The Ecne analysis is complementary to the formal security analysis.

Theorems in Section~\ref{sec:security} rely on the specified
cryptographic assumptions and on the deterministic correctness of the
deployed contract and circuit.

Ecne does not replace those assumptions and is not included as a term in
the reduction.

Its role is narrower: it provides additional implementation-level
evidence that the R1CS does not contain an obvious under-constrained
signal under the stated analysis methodology.

\begingroup
\footnotesize
\setlength{\itemsep}{0pt}
\begin{thebibliography}{99}

\bibitem{groth16} Jens Groth. \emph{On the Size of Pairing-Based
Non-Interactive Arguments.} In \emph{Advances in Cryptology ---
EUROCRYPT 2016}, LNCS 9666, pp.\ 305--326. Springer, 2016. IACR ePrint
2016/260.

\bibitem{poseidon} Lorenzo Grassi, Dmitry Khovratovich, Christian
Rechberger, Arnab Roy, and Markus Schofnegger. \emph{Poseidon: A New
Hash Function for Zero-Knowledge Proof Systems.} In \emph{USENIX
Security Symposium}, 2021.

\bibitem{rayls} M. Yaksetig and J. Xu. \emph{Rayls: A Novel Design for
CBDCs.} IACR ePrint 2025/1639, 2025.

\bibitem{rayls2} M. Yaksetig, P.\,M.\,F. Pereira, S. Yang, M. Nejadgholi,
and J. Xu. \emph{Rayls II: Fast, Private, and Compliant CBDCs.} IACR
ePrint 2025/1638, 2025.

\bibitem{zeto} LF Decentralized Trust, Paladin Project. \emph{Zeto:
Privacy-preserving implementations of fungible and non-fungible tokens,
using UTXO as the underlying transaction model.} GitHub repository,
accessed September 2026.
\url{https://github.com/LF-Decentralized-Trust-labs/paladin}.

\bibitem{zkp2p} ZKP2P / Peer. \emph{Peer Documentation.} Documentation
portal, accessed September 2026.

\bibitem{canton} Digital Asset. \emph{Canton Network: A Network of
Networks for Smart Contract Applications.} White paper, updated January
2024.

\bibitem{deco} F. Zhang, D. Maram, H. Malvai, S. Goldfeder, and A.
Juels. \emph{DECO: Liberating Web Data Using Decentralized Oracles for
TLS.} In \emph{Proceedings of the ACM Conference on Computer and
Communications Security (CCS)}, 2020.

\bibitem{mlkem} National Institute of Standards and Technology.
\emph{FIPS 203: Module-Lattice-Based Key-Encapsulation Mechanism
Standard.} August 2024.
\url{https://csrc.nist.gov/pubs/fips/203/final}.

\bibitem{qurrency} A.\,R. Choudhuri, S. Garg, M. Gregoire, K. Lee, M.
Lodder, H. Montgomery, G.\,V. Policharla, and J. Zhang. \emph{Qurrency: A
Quantum-Secure, Private, and Auditable Platform for Digital Assets.}
IACR ePrint 2026/015, 2026.

\bibitem{ecne} F. Wang. \emph{Ecne: An Engine for Verifying the
Soundness of R1CS Constraints.} 0xPARC, 2022.
\url{https://github.com/franklynwang/EcneProject}.

\bibitem{qed2} S. Pailoor, Y. Chen, F. Wang, C. Rodr\'iguez, J. Van
Gaffen, J. Morton, M. Chu, B. Gu, Y. Feng, and I. Dillig. \emph{Automated
Detection of Underconstrained Circuits for Zero-Knowledge Proofs.} In
\emph{Proceedings of the ACM SIGPLAN Conference on Programming Language
Design and Implementation (PLDI)}, 2023. IACR ePrint 2023/512.

\bibitem{circomspect} Trail of Bits. \emph{Circomspect: A Static
Analyzer and Linter for the Circom Zero-Knowledge DSL.} 2022.
\url{https://github.com/trailofbits/circomspect}.

\bibitem{spicatalog} Banco Central do Brasil. \emph{Cat\'alogo de
Servi\c{c}os do SFN, Volume VI --- Pagamentos Instant\^aneos}, vers\~ao
5.12.

\bibitem{cnpjalfa} Receita Federal do Brasil. \emph{Instru\c{c}\~ao
Normativa RFB n\textsuperscript{o} 2.229, de 15 de outubro de 2024.}

\end{thebibliography}
\endgroup

\end{document}
